\documentclass[11pt]{article}
\usepackage{enumitem}
\usepackage{listings}
\lstset{basicstyle=\ttfamily, breaklines=true}
\begin{document}
  \begin{enumerate}
    \item [2.1][5]$<$§2.2$>$ For the following C statement, write the corresponding RISC-V
assembly code. Assume that the C variables f, g, and h, have already been placed
in registers x5, x6, and x7 respectively. Use a minimal number of RISC-V assembly
instructions.
\begin{lstlisting}
  f = g + (h - 5);
\end{lstlisting}
\begin{lstlisting}
  addi x7, x7, -5
  add x5, x6, x7
\end{lstlisting}
\item [2.2][5] $<$§2.2$>$ Write a single C statement that corresponds to the two RISC-V
assembly instructions below.
\begin{lstlisting}
  add f, g, h
  add f, i, f

  //Solution
  f = i + g + h
\end{lstlisting}
\item [2.3][5] $<$§§2.2, 2.3$>$ For the following C statement, write the corresponding
RISC-V assembly code. Assume that the variables f, g, h, i, and j are assigned to
registers x5, x6, x7, x28, and x29, respectively. Assume that the base address
of the arrays A and B are in registers x10 and x11, respectively.
\begin{lstlisting}
  B[8] = A[i-j];

  sub x28, x28, x29
  slli x28, x28, 8
  add x10, x10, x28

  ld x5, 0(x10)
  sd x5, 64(x11)
\end{lstlisting}
\item [2.4][10] $<$§§2.2, 2.3$>$ For the RISC-V assembly instructions below, what is the
corresponding C statement? Assume that the variables f, g, h, i, and j are assigned
to registers x5, x6, x7, x28, and x29, respectively. Assume that the base
address of the arrays A and B are in registers x10 and x11, respectively.
\begin{lstlisting}
  slli x30, x5, 3      // x30 = f*8
  add  x30, x10, x30   // x30 = &A[f]
  slli x31, x6, 3      // x31 = g*8
  add  x31, x11, x31   // x31 = &B[g]
  ld   x5, 0(x30)      //   f = A[f]

  addi x12, x30 , 8    //x12 = A[f] + 8
  ld   x30, 0(x12)     //A[f] = A[f+1]
  add  x30, x30, x5    //A[f+1] = A[f+1] + A[f]
  sd   x30, 0(x31)     //B[g] = A[f+1] + A[f]

  //Solution
  B[g] = A[f+1] + A[f];
\end{lstlisting}
\item [2.5][5] $<$§2.3$>$ Show how the value \lstinline|0xabcdef12| would be arranged in memory
of a little-endian and a big-endian machine. Assume the data are stored starting at
address 0 and that the word size is 4 bytes.
\begin{lstlisting}
  little-endian:0x12efcdab
  big-endian:0xabcdef12
\end{lstlisting}
\item[2.6][5] $<$§2.4$>$ Translate \lstinline|0xabcdef12| into decimal.
\item[2.7] [5] $<$§§2.2, 2.3$>$ Translate the following C code to RISC-V. Assume that the
variables f, g, h, i, and j are assigned to registers x5, x6, x7, x28, and x29,
respectively. Assume that the base address of the arrays A and B are in registers x10
and x11, respectively. Assume that the elements of the arrays A and B are 8-byte
words:
\begin{lstlisting}
  B[8] = A[i] + A[j]; 

  //Solution. Use x15 and x16 as temporary registers.
  slli x30, x29, 3      // x30 = j*8
  add x30, x10, x30    // x30 = &A[j]
  ld x15, 0(x30)        // x15 = A[j]
  slli x31, x28, 3      // x31 = i*8
  add x31, x10, x31    // x31 = &A[i]
  ld x16, 0(x31)        // x16 = A[i]
  add x16, x16, x15
  sd x16, 64(x11)
\end{lstlisting}
\item[2.8][10] $<$§§2.2, 2.3$>$ Translate the following RISC-V code to C. Assume that the
variables f, g, h, i, and j are assigned to registers x5, x6, x7, x28, and x29, respectively. Assume that the base address of the arrays A and B are in registers x10
and x11, respectively.
  \begin{lstlisting}
    addi x30, x10, 8 //x30 = x10 + 8= A[1]
    addi x31, x10, 0  //x31 = x10 + 0 = A[0]
    sd x31, 0(x30)    //A[1] = A[0]
    ld x30, 0(x30)    // A[1] = A[0]
    add x5, x30, x31  // x5 = A[1]+A[0]

    //Solution
    f = A[1]+A[0]
  \end{lstlisting}
  \end{enumerate}
\end{document}